\section{НАУЧНАЯ НОВИЗНА}

В банковской сфере наблюдается устойчивый тренд перехода от использования типовых готовых программных продуктов к разработке собственных интеллектуальных систем сопровождения корпоративных хранилищ данных (КХД). Так, например, ПАО «Промсвязьбанк» (ПСБ) инициировал сотрудничество с профильными образовательными учреждениями, в частности с Московским авиационным институтом, где была создана продуктовая модель для анализа кредитоспособности клиентов. Это демонстрирует растущую потребность банков в кастомизированных, гибких и интеллектуальных решениях, способных адаптироваться к специфике внутренних бизнес-процессов и требованиям регуляторов.

Современные инструменты интеграции и сопровождения данных, такие как Oracle Data Integrator и Oracle GoldenGate, широко применяются для решения задач интеграции, репликации и управления потоками данных. Несмотря на широкое применение Oracle Data Integrator и GoldenGate для интеграции и репликации данных, текущие процессы зачастую требуют ручной настройки «запинок» — специалист формирует, из каких систем и таблиц извлекать данные и как организовать их трансформацию. Это ограничивает скорость адаптации и гибкость систем, требует высокой квалификации и значительных временных затрат.

Традиционные алгоритмы оптимизации SQL-запросов в этих решениях не всегда справляются с динамически меняющимися сценариями, что негативно влияет на производительность.

Интеллектуальные агенты на базе больших языковых моделей способны значительно расширить возможности сопровождения КХД. Они могут:
\begin{itemize}
    \item самостоятельно анализировать структуру и содержимое данных;
    \item автоматически выявлять релевантные таблицы и источники информации;
    \item формировать и оптимизировать SQL-запросы с учётом текущей нагрузки и особенностей базы данных;
    \item обеспечивать полноценный цикл сопровождения с помощью методов обратного проектирования (reverse engineering), минимизируя ручное вмешательство и снижая риски ошибок.
\end{itemize}


Такой подход открывает новые горизонты для повышения эффективности, устойчивости и автоматизации корпоративных хранилищ данных в банковской сфере.

Данный подход требует высокой квалификации, занимает значительное время и ограничивает гибкость и адаптивность систем. Кроме того, традиционные алгоритмы оптимизации SQL-запросов, применяемые в этих инструментах, не всегда способны эффективно справляться с динамически меняющимися, комплексными сценариями обработки данных, что сказывается на производительности и времени отклика.

Научная новизна данной работы заключается в применении больших языковых моделей (LLM) для построения интеллектуальной системы сопровождения КХД, которая обеспечивает:
\begin{itemize}
    \item автоматизированный и адаптивный анализ SQL-запросов и процедур с динамическим обучением и самосовершенствованием алгоритмов, направленный на повышение производительности и снижение нагрузки на инфраструктуру;
    \item интеллектуальное выявление аномалий и диагностику ошибок на основе комплексного анализа логов, метрик и структурированных данных с генерацией понятных и обоснованных рекомендаций на естественном языке, что облегчает взаимодействие с техническими специалистами;
    \item автоматическую генерацию регламентированных отчетов, полностью соответствующих требованиям банковского регулирования, с возможностью оперативного обновления моделей в соответствии с изменениями нормативной базы;
    \item интеграцию ИИ-агентов в процессы мониторинга и сопровождения КХД, обеспечивающую проактивный контроль, предиктивную аналитику и значительное сокращение трудозатрат на сопровождение.
\end{itemize}

Особенно перспективным является использование методов обратного проектирования (reverse engineering), благодаря которым ИИ-агенты способны полностью автоматизировать цикл сопровождения — от анализа архитектуры данных и оптимизации процессов интеграции до формирования отчетности. Это выводит разработанную систему за рамки традиционных инструментов, создавая комплексный, интеллектуальный и адаптивный продукт, отвечающий современным вызовам банковского сектора.
